% Style based on the supplied Waseem_Generic_AI_ML_GenAI_AgenticAI.tex.
% Python replaces template tokens. Edit layout here, content in JSON/portfolio JS.
\documentclass[9pt,letterpaper]{extarticle}
\usepackage[left=0.69in,right=0.69in,top=0.45in,bottom=0.45in]{geometry}
\usepackage[T1]{fontenc}
\input{glyphtounicode}
\pdfgentounicode=1
\usepackage[utf8]{inputenc}
\usepackage{lmodern}
\usepackage{microtype}
\usepackage{pifont}
\usepackage{xcolor}
\usepackage{enumitem}
\usepackage{titlesec}
\usepackage[hidelinks]{hyperref}
\usepackage{tabularx,array,parskip,ragged2e,needspace}
\pagestyle{empty}
\setlength{\parindent}{0pt}
\setlength{\parskip}{1pt}
\definecolor{accent}{HTML}{155B8A}
\titleformat{\section}{\large\bfseries\color{accent}}{}{0pt}{}[\vspace{0.5pt}\color{accent}\titlerule]
\titlespacing*{\section}{0pt}{3pt}{2pt}
% Compact subsections beneath the main accent-rule sections.
\titleformat{\subsection}{\normalsize\bfseries\color{accent}}{}{0pt}{}
\titlespacing*{\subsection}{0pt}{3pt}{1pt}
\setlist[itemize]{leftmargin=1.15em,itemsep=0.8pt,topsep=1pt,parsep=0pt,partopsep=0pt}
\newcommand{\role}[4]{%
  \Needspace{4\baselineskip}\textbf{#1}\hfill\textbf{#2}\\[-1pt]
  \textit{#3}\hfill\textit{#4}\par}
\newcommand{\entry}[2]{\Needspace{4\baselineskip}\textbf{#1}\hfill\textit{#2}\par}
\newcommand{\skillrow}[2]{\textbf{#1:} & #2\\}
\hypersetup{pdftitle={Waseem Raza, PhD - Reusable Academic CV \textbar{} ECE, Communications, Sensing and Applied AI \textbar{} 2027},pdfauthor={Waseem Raza, PhD}}
\begin{document}
\hyphenpenalty=10000\exhyphenpenalty=10000\pretolerance=10000\emergencystretch=2em\widowpenalty=10000\clubpenalty=10000
\begin{center}
{\LARGE\bfseries Waseem Raza, PhD}{\normalsize\mdseries \enspace (Electrical and Computer Engineering)}\\[2pt]
{\large\bfseries Electrical and Computer Engineering \textbar{} Research, Teaching \& Applied AI}\\[1pt]
{\small Wireless Systems • Signal Processing • Machine Learning • Intelligent Sensing}\\[3pt]
{\footnotesize \ding{41}\enspace \href{mailto:waseem.placements@gmail.com}{waseem.placements@gmail.com} \textbar{} \href{mailto:waseemraza844@gmail.com}{waseemraza844@gmail.com}\\[1pt]New Jersey, United States \quad \ding{37}\enspace \href{tel:+17328039703}{+1 732-803-9703}\\[1pt]\href{https://www.linkedin.com/in/waseem-raza-phd-9a737858/}{LinkedIn} \textbar{} \href{https://github.com/WaseemRaza844}{GitHub} \textbar{} \href{https://scholar.google.com/citations?user=d58Drr0AAAAJ\&hl=en}{Google Scholar} \textbar{} \href{https://waseemraza844.github.io/portfolio/learning.html}{Coursera Courses} \textbar{} \href{https://waseemraza844.github.io/portfolio/}{Portfolio}}
\end{center}
\vspace{-5pt}
\Needspace{3\baselineskip}
\section{Academic Profile}
PhD in Electrical and Computer Engineering (University of Oklahoma, 2025) with undergraduate teaching in Computer Programming, Communication Systems and Wireless Communication. Research on reliable AI/ML for wireless networks includes propagation prediction, generative augmentation, transfer learning, outage detection and network self-healing. Related work in wearable sensing and time-series modeling. Engineering R\&D at AT\&T Labs, Futurewei and AT\&T, alongside student mentoring and communications laboratory instruction.
\Needspace{2\baselineskip}
\subsection{Research Areas}
AI/ML for Wireless and Networked Systems; Reliable and Data-Efficient Learning; Propagation Modeling and Digital Twins; Signal Processing and Intelligent Sensing; Network Resilience and Optimization; Wearable Sensors and Time-Series Analytics\par
\Needspace{3\baselineskip}
\section{Education}
\entry{PhD, Electrical and Computer Engineering}{09/2019 - 07/2025}
University of Oklahoma, Norman, Oklahoma, USA\par
\entry{MSc, Telecommunication Engineering}{08/2014 - 08/2016}
University of Engineering and Technology, Taxila, Pakistan\par
\entry{BSc, Telecommunication Engineering}{11/2010 - 06/2014}
University of Engineering and Technology, Taxila, Pakistan\par
\Needspace{2\baselineskip}
\subsection{Doctoral Research}
Dissertation on AI/ML for wireless networks, covering robust propagation and digital-twin modeling, domain adaptation, generative data augmentation, and machine-learning-assisted outage detection and network self-healing.\par
\Needspace{3\baselineskip}
\section{Teaching and Mentoring}
\Needspace{6\baselineskip}
\Needspace{2\baselineskip}
\subsection{University-Level Teaching}
\Needspace{6\baselineskip}
\role{Lecturer, Department of Technology}{12/2017 - 09/2019}{University of Lahore, Pakistan}{Lahore, Pakistan}
\begin{itemize}\interlinepenalty=10000
\item Taught programming, communication systems, and wireless communication with practical examples.
\item Advised senior-design students on projects, academic plans, and early careers.
\item Aligned course materials with learning outcomes and accreditation standards.
\item Supported accreditation, exams, technical workshops, and student activities.
\end{itemize}
\vspace{1pt}
{\small\textbf{Teaching areas:} Computer Programming (previously taught); Communication Systems (previously taught); Wireless Communication (previously taught); Signals and Digital Signal Processing (technical preparation); Computer Networks and Applied Machine Learning (technical preparation); Engineering Simulation and Laboratory Instruction}\par
\Needspace{2\baselineskip}\noindent\textbf{Student mentoring and research development}\par
\begin{itemize}\interlinepenalty=10000
\item Advised undergraduate student projects and supported Bachelor's and Master's research/thesis projects in engineering software, simulations and communications systems.
\item Helped students scope experiments, analyze results and communicate limitations in project reports and technical presentations.
\item Supported hands-on communications laboratories, programming exercises, assignment preparation and assessment.
\end{itemize}
\Needspace{3\baselineskip}
\section{Technical Skills, Engineering Methods \& Research Tools}
\Needspace{2\baselineskip}\noindent{\small\textbf{Programming \& Software:} Python, C/C++, MATLAB, SQL; debugging, Git/GitHub, Linux, Jupyter and reproducible workflows}\par
\Needspace{2\baselineskip}\noindent{\small\textbf{Machine Learning \& Analytics:} PyTorch, TensorFlow, Keras, scikit-learn and XGBoost; generative and transfer learning, reinforcement learning, classification and evaluation}\par
\Needspace{2\baselineskip}\noindent{\small\textbf{Signal Processing \& Networking:} Communication systems, signal analysis, propagation, RF and cellular-network measurements, LTE/4G, 5G-NR and network optimization}\par
\Needspace{2\baselineskip}\noindent{\small\textbf{Simulation \& Data:} NumPy, SciPy, Pandas, Matplotlib, simulation, statistical validation, distribution-shift testing and technical reporting}\par
\Needspace{2\baselineskip}\noindent{\small\textbf{Teaching \& Mentorship:} Undergraduate Technology instruction, programming exercises, communications labs, student projects, academic assessment and feedback}\par
\Needspace{3\baselineskip}
\section{Academic and Industry Research Experience}
\Needspace{6\baselineskip}
\Needspace{2\baselineskip}
\subsection{University Research Appointments}
\Needspace{6\baselineskip}
\role{Graduate Research Assistant}{09/2019 - 07/2025}{AI4Networks, University of Oklahoma}{Tulsa, OK}
\begin{itemize}\interlinepenalty=10000
\item Built 3GPP-compliant Python simulation tools and an AI-enabled 5G research testbed.
\item Built GAN, transfer-learning, autoencoder, and reinforcement-learning models for wireless data.
\item Led 100+ experiments testing wireless ML under domain shift and limited labeled data.
\item Developed self-healing and outage-management methods, improving resilience by 30\%.
\item Published wireless AI research and mentored junior researchers in experimental design.
\end{itemize}
\vspace{1pt}
\Needspace{6\baselineskip}
\role{Research Associate, ComSens Lab}{06/2016 - 12/2017}{COMSATS University Islamabad, Pakistan}{Islamabad, Pakistan}
\begin{itemize}\interlinepenalty=10000
\item Modeled wireless communication systems through simulation, analysis, and technical reports.
\item Mentored Bachelor's and Master's theses in experiment design and research presentation.
\item Reviewed student reports and manuscripts for rigor and publication readiness.
\item Supported proposals, literature reviews, lab seminars, and research coordination.
\end{itemize}
\vspace{1pt}
\Needspace{6\baselineskip}
\role{Full-Time Research Scholar, Communication Systems Lab}{08/2014 - 02/2016}{University of Engineering and Technology, Taxila, Pakistan}{Taxila, Pakistan}
\begin{itemize}\interlinepenalty=10000
\item Researched wireless communication systems while completing an MSc in Telecommunications.
\item Supported communications courses, labs, assignments, grading, and examinations.
\item Led wireless communication labs and prepared technical manuals and experiments.
\item Developed course and student-project materials for practical wireless engineering.
\end{itemize}
\vspace{1pt}
\Needspace{6\baselineskip}
\Needspace{2\baselineskip}
\subsection{Industrial Research and Engineering}
\Needspace{6\baselineskip}
\role{Machine Learning Engineer}{02/2025 - Present}{AT\&T, contractor via Innovatix Inc.}{Bedminster, NJ}
\begin{itemize}\interlinepenalty=10000
\item Applied domain-adaptive ML to telecom service assurance and model validation.
\item Designed cross-domain evaluation and error-analysis methods for operational ML.
\item Worked with engineering and public-safety teams on reliable network assurance.
\end{itemize}
\vspace{1pt}
\Needspace{6\baselineskip}
\role{Data Scientist and AI/ML Engineer}{09/2024 - 12/2024}{Impac Exploration Services}{Weatherford, OK}
\begin{itemize}\interlinepenalty=10000
\item Built LLM-enabled analytics to interpret geochemical well data for decision support.
\item Built Python pipelines in Azure Databricks, Functions, and Synapse Analytics.
\item Automated exploration-data ETL, cutting preparation time by 40\%.
\item Used data-validation and quality controls to improve output reliability by 20\%.
\item Turned geologists' technical needs into structured datasets and useful analytics.
\end{itemize}
\vspace{1pt}
\Needspace{6\baselineskip}
\role{Advanced-Degrees Intern}{06/2024 - 08/2024}{AT\&T Labs}{Austin, TX}
\begin{itemize}\interlinepenalty=10000
\item Analyzed 500K+ LTE/5G-NR logs with ML, reducing diagnostic time by 20\%.
\item Improved 5G NR-NTN cell selection by 15\% through data-driven validation.
\item Presented operational data insights to next-generation wireless engineering teams.
\end{itemize}
\vspace{1pt}
\Needspace{6\baselineskip}
\role{Advanced Wireless Technology Development Intern}{09/2022 - 01/2023}{Futurewei Technologies}{Schaumburg, IL}
\begin{itemize}\interlinepenalty=10000
\item Improved wireless CSI compression by 25\% with transfer learning and autoencoders.
\item Ran 50+ experiments on wireless ML robustness under changing data conditions.
\item Translated research results into recommendations for future wireless networks.
\end{itemize}
\vspace{1pt}
\Needspace{6\baselineskip}
\role{Advanced-Degrees Intern}{06/2022 - 08/2022}{AT\&T Labs}{Bedminster, NJ}
\begin{itemize}\interlinepenalty=10000
\item Improved 5G standalone user-timing precision by 20\% with ML and statistics.
\item Analyzed user traces with SQL and presented results to engineering leadership.
\item Turned operational traces into indicators for network diagnosis and optimization.
\end{itemize}
\vspace{1pt}
\Needspace{3\baselineskip}
\section{Selected Peer-Reviewed Publications}
\Needspace{4\baselineskip}\noindent\begin{tabularx}{\linewidth}{@{}>{\raggedright\arraybackslash}p{0.68\linewidth}@{\hspace{8pt}}>{\raggedleft\arraybackslash}X@{}}
\textbf{Real-World Medication Adherence Monitoring Using Hybrid Neural Networks and Smartwatches} & {\scriptsize\textit{HealthINF \textbar{} Accepted}\newline 2026}\\
\end{tabularx}\vspace{1pt}
\Needspace{4\baselineskip}\noindent\begin{tabularx}{\linewidth}{@{}>{\raggedright\arraybackslash}p{0.68\linewidth}@{\hspace{8pt}}>{\raggedleft\arraybackslash}X@{}}
\textbf{\href{https://doi.org/10.1016/j.comcom.2025.108129}{AI-Powered Resilience: A Dual-Approach for Outage Management in Dense Cellular Networks}} & {\scriptsize\textit{Computer Communications}\newline 2025 \textbar{} \href{https://doi.org/10.1016/j.comcom.2025.108129}{DOI}}\\
\end{tabularx}\vspace{1pt}
\Needspace{4\baselineskip}\noindent\begin{tabularx}{\linewidth}{@{}>{\raggedright\arraybackslash}p{0.68\linewidth}@{\hspace{8pt}}>{\raggedleft\arraybackslash}X@{}}
\textbf{\href{https://doi.org/10.3390/s25247651}{On Multi-Parameter Optimization and Proactive Reliability in 5G and Beyond Cellular Networks}} & {\scriptsize\textit{Sensors 2025}\newline 2025 \textbar{} \href{https://doi.org/10.3390/s25247651}{DOI}}\\
\end{tabularx}\vspace{1pt}
\Needspace{4\baselineskip}\noindent\begin{tabularx}{\linewidth}{@{}>{\raggedright\arraybackslash}p{0.68\linewidth}@{\hspace{8pt}}>{\raggedleft\arraybackslash}X@{}}
\textbf{\href{https://doi.org/10.1109/PIMRC59610.2024.10817403}{An AI-Driven Framework for Enhancing Resilience in Propagation Models to Enable Digital Twin}} & {\scriptsize\textit{IEEE PIMRC}\newline 2024 \textbar{} \href{https://doi.org/10.1109/PIMRC59610.2024.10817403}{DOI}}\\
\end{tabularx}\vspace{1pt}
\Needspace{4\baselineskip}\noindent\begin{tabularx}{\linewidth}{@{}>{\raggedright\arraybackslash}p{0.68\linewidth}@{\hspace{8pt}}>{\raggedleft\arraybackslash}X@{}}
\textbf{\href{https://www.ai4networks.com/files/conferences/c-24-7.pdf}{Refining Wireless Propagation Models using Domain-Informed GANs amid Data Scarcity}} & {\scriptsize\textit{IEEE VTC-Fall 2024}\newline 2024 \textbar{} \href{https://www.ai4networks.com/files/conferences/c-24-7.pdf}{Link}}\\
\end{tabularx}\vspace{1pt}
\Needspace{4\baselineskip}\noindent\begin{tabularx}{\linewidth}{@{}>{\raggedright\arraybackslash}p{0.68\linewidth}@{\hspace{8pt}}>{\raggedleft\arraybackslash}X@{}}
\textbf{Data-Driven Intelligent Outage Management for High Shadowing Environments in 5G\&B Networks} & {\scriptsize\textit{IEEE SmartNets}\newline 2024}\\
\end{tabularx}\vspace{1pt}
\Needspace{4\baselineskip}\noindent\begin{tabularx}{\linewidth}{@{}>{\raggedright\arraybackslash}p{0.68\linewidth}@{\hspace{8pt}}>{\raggedleft\arraybackslash}X@{}}
\textbf{\href{https://doi.org/10.1109/MNET.006.2100377}{An AI-Enabled Framework to Defend Ingenious MDT-Based Attacks on the Emerging Zero Touch Cellular Networks}} & {\scriptsize\textit{IEEE Network, 38(1), 228-237}\newline 2024 \textbar{} \href{https://doi.org/10.1109/MNET.006.2100377}{DOI}}\\
\end{tabularx}\vspace{1pt}
\Needspace{4\baselineskip}\noindent\begin{tabularx}{\linewidth}{@{}>{\raggedright\arraybackslash}p{0.68\linewidth}@{\hspace{8pt}}>{\raggedleft\arraybackslash}X@{}}
\textbf{\href{https://arxiv.org/abs/2402.15795}{Positioning Error Impact Compensation through Data-Driven Optimization in User-Centric Networks}} & {\scriptsize\textit{IEEE GLOBECOM 2023}\newline 2023 \textbar{} \href{https://arxiv.org/abs/2402.15795}{Link}}\\
\end{tabularx}\vspace{1pt}
\Needspace{4\baselineskip}\noindent\begin{tabularx}{\linewidth}{@{}>{\raggedright\arraybackslash}p{0.68\linewidth}@{\hspace{8pt}}>{\raggedleft\arraybackslash}X@{}}
\textbf{\href{https://arxiv.org/abs/2203.15899}{Toward a Hybrid RF/Optical Lunar Communication System (LunarComm)}} & {\scriptsize\textit{IEEE Network}\newline 2022 \textbar{} \href{https://arxiv.org/abs/2203.15899}{Link}}\\
\end{tabularx}\vspace{1pt}
\Needspace{4\baselineskip}\noindent\begin{tabularx}{\linewidth}{@{}>{\raggedright\arraybackslash}p{0.68\linewidth}@{\hspace{8pt}}>{\raggedleft\arraybackslash}X@{}}
\textbf{Towards Positioning Error Impact Characterization and Minimization in User-Centric RAN} & {\scriptsize\textit{IEEE WCNC}\newline 2022}\\
\end{tabularx}\vspace{1pt}
\Needspace{4\baselineskip}\noindent\begin{tabularx}{\linewidth}{@{}>{\raggedright\arraybackslash}p{0.68\linewidth}@{\hspace{8pt}}>{\raggedleft\arraybackslash}X@{}}
\textbf{\href{https://doi.org/10.1016/j.aeue.2020.153244}{Unification of RF Energy Harvesting Schemes under Mixed Rayleigh-Rician Fading Channels}} & {\scriptsize\textit{AEÜ - International Journal of Electronics and Communications}\newline 2020 \textbar{} \href{https://doi.org/10.1016/j.aeue.2020.153244}{DOI}}\\
\end{tabularx}\vspace{1pt}
\Needspace{3\baselineskip}
\section{Selected Research and Engineering Platforms}
\entry{SyntheticNET: 3GPP-Compliant Network Simulation}{Research Project \textbar{} 2020-2024}
Developed a Python-based, 3GPP-aligned simulation environment for realistic mobility, propagation, and cellular-network experimentation when measured data is sparse or costly. Supported controlled data generation and repeatable studies in propagation modeling, network optimization, digital twins, and resilient wireless AI.\par
{\small\textit{Tools \& Methods:} Python, 3GPP, Wireless Simulation, Mobility Modeling, Propagation Modeling, Synthetic Data, Experiment Design, Digital Twins, Network Analytics\par}
\vspace{1pt}
\entry{TurboRAN Experimental Platform}{Research Infrastructure \textbar{} 2020-2024}
Built and used a multi-band, multi-tier experimental platform for evaluating AI-enabled radio-access-network methods under controlled network conditions. Supported repeatable measurement collection, testbed experimentation, RAN optimization studies, and validation of emerging wireless-network algorithms.\par
{\small\textit{Tools \& Methods:} 5G Testbeds, RAN, Network Measurement, Experiment Design, Wireless Prototyping, Multi-Band Systems, Performance Evaluation, RAN Optimization\par}
\vspace{1pt}
\entry{Propagation Models \& Digital Twins}{Research Project \textbar{} 2024}
Developed robust data-driven radio-propagation models for wireless digital twins using domain-informed generative learning and scarce measured data. Combined 3GPP-informed pretraining, GAN-based augmentation, domain adaptation, and statistical validation to improve resilience under distribution shift.\par
{\small\textit{Tools \& Methods:} Propagation Modeling, Digital Twins, GANs, Domain Adaptation, 3GPP Models, Synthetic Data, Statistical Validation, Wireless AI, Distribution Shift\par}
\vspace{1pt}
\entry{Wearable Intelligence for Health Monitoring}{Research Project \textbar{} 2024-2026}
Developed deep-learning workflows for recognizing medication-taking and dietary activities from smartwatch accelerometer and gyroscope signals in free-living settings. Evaluated windowing, feature extraction, CNN/LSTM-style sequence models, cross-user generalization, and recall-oriented performance for health-monitoring applications.\par
{\small\textit{Tools \& Methods:} Wearable AI, Smartwatch IMU, Time Series, Deep Learning, CNN/LSTM, Feature Engineering, Cross-User Evaluation, Activity Recognition, Healthcare AI\par}
\vspace{1pt}
\entry{AI for Telecom Network \& Security Optimization}{Certification Project \textbar{} 2026}
Applied AI concepts to telecommunications network optimization, security analysis, and customer-experience improvement through a structured specialization project. Framed telecom use cases around operational data, AI-assisted analysis, optimization opportunities, and practical service-improvement recommendations.\par
{\small\textit{Tools \& Methods:} AI for Telecommunications, Network Optimization, Telecom Security, Customer Experience, Python, TensorFlow, Operational Analytics, Service Improvement\par}
\vspace{1pt}
\Needspace{3\baselineskip}
\section{Academic Service and Honors}
\Needspace{2\baselineskip}
\subsection{Selected Academic Service}
\begin{itemize}\interlinepenalty=10000
\item Contributed to engineering course material, student evaluation, examinations and curriculum-related activities.
\item Supported laboratory technical documentation, research seminars and student projects.
\end{itemize}
\Needspace{2\baselineskip}
\subsection{Selected Honors}
\entry{William H. Barkow Scholarship}{2024}
{\small University of Oklahoma, Electrical and Computer Engineering\par}
\entry{Clyde L. Farrar Endowment Scholarship}{2024}
{\small University of Oklahoma\par}
\entry{P.H. Robinson Fellowship}{2023}
{\small University of Oklahoma\par}
\entry{FAS-Charles Forgan Scholarship}{2023}
{\small University of Oklahoma\par}
\entry{Nettie Vincent Boggs Graduate Endowment}{2022}
{\small University of Oklahoma\par}
\entry{SHS-Paul D. Massad Scholarship}{2022}
{\small University of Oklahoma\par}
\entry{Cleo Cross Memorial Scholarship}{2021}
{\small University of Oklahoma\par}
\entry{Marjorie Haskell Engineering Scholarship}{2021}
{\small University of Oklahoma\par}
\entry{Fully Funded PhD / Graduate Research Assistantship}{2019}
{\small University of Oklahoma\par}
\Needspace{3\baselineskip}
\section{Selected Professional Development}
\Needspace{4\baselineskip}\noindent\begin{tabularx}{\linewidth}{@{}>{\raggedright\arraybackslash}p{0.59\linewidth}@{\hspace{7pt}}>{\raggedleft\arraybackslash}X@{}}
\textbf{1. \href{https://www.coursera.org/account/accomplishments/verify/VU3Y5TBLVLZT}{Machine Learning}} & {\scriptsize \href{https://www.coursera.org/account/accomplishments/verify/VU3Y5TBLVLZT}{Stanford/Coursera} \textbar{} Completed \textbar{} Mar 2020}\\
\end{tabularx}\vspace{-2pt}
\noindent{\footnotesize Supervised and unsupervised learning, regularization, model evaluation, and practical predictive modeling.}\par\vspace{2pt}
\Needspace{4\baselineskip}\noindent\begin{tabularx}{\linewidth}{@{}>{\raggedright\arraybackslash}p{0.59\linewidth}@{\hspace{7pt}}>{\raggedleft\arraybackslash}X@{}}
\textbf{2. \href{https://waseemraza844.github.io/portfolio/certificates/cloud-devops/C7. Coursera WRA89JEJU8EE.pdf}{Python for Data Science, AI \& Development}} & {\scriptsize \href{https://www.coursera.org/account/accomplishments/verify/WRA89JEJU8EE}{IBM/Coursera} \textbar{} Completed \textbar{} Jul 2026}\\
\end{tabularx}\vspace{-2pt}
\noindent{\footnotesize Python programming, data analysis, APIs, and reproducible data-science workflows.}\par\vspace{2pt}
\Needspace{4\baselineskip}\noindent\begin{tabularx}{\linewidth}{@{}>{\raggedright\arraybackslash}p{0.59\linewidth}@{\hspace{7pt}}>{\raggedleft\arraybackslash}X@{}}
\textbf{3. \href{https://www.coursera.org/account/accomplishments/verify/Y6NP4GWLS4DZ}{Introduction to Deep Learning \& Neural Networks with Keras}} & {\scriptsize \href{https://www.coursera.org/account/accomplishments/verify/Y6NP4GWLS4DZ}{IBM/Coursera} \textbar{} Completed \textbar{} Aug 2026}\\
\end{tabularx}\vspace{-2pt}
\noindent{\footnotesize Neural networks, training workflows, optimization, loss functions, and model evaluation with Keras.}\par\vspace{2pt}
\Needspace{4\baselineskip}\noindent\begin{tabularx}{\linewidth}{@{}>{\raggedright\arraybackslash}p{0.59\linewidth}@{\hspace{7pt}}>{\raggedleft\arraybackslash}X@{}}
\textbf{4. \href{https://waseemraza844.github.io/portfolio/certificates/wireless/4g-network-fundamentals.pdf}{4G Network Fundamentals}} & {\scriptsize \href{https://www.coursera.org/account/accomplishments/verify/DLW1YN474VZJ}{Mines-Télécom/Coursera} \textbar{} Completed \textbar{} Aug 2026}\\
\end{tabularx}\vspace{-2pt}
\noindent{\footnotesize LTE radio and core networks, cellular architecture, mobility management, and 4G fundamentals.}\par\vspace{2pt}
\Needspace{4\baselineskip}\noindent\begin{tabularx}{\linewidth}{@{}>{\raggedright\arraybackslash}p{0.59\linewidth}@{\hspace{7pt}}>{\raggedleft\arraybackslash}X@{}}
\textbf{5. \href{https://waseemraza844.github.io/portfolio/certificates/wireless/c1-ai-found-indus-overview.pdf}{AI Foundations \& Industry Overview for Telecommunication}} & {\scriptsize \href{https://www.coursera.org/account/accomplishments/verify/3900EVOH4I1Y}{Coursera} \textbar{} Completed \textbar{} Jun 2026}\\
\end{tabularx}\vspace{-2pt}
\noindent{\footnotesize AI applications in telecom networks, operations analytics, service assurance, and automation.}\par\vspace{2pt}
\end{document}
